\section{课程定位与主线}

本讲由 UT Austin 的 Peter Stone 主讲，题目是 Autonomous Agents。它并不是在讲一套新的 LLM 框架，而是在把 ``Agent'' 这个概念重新放回 AI 长期研究脉络：\textbf{embodiment、interaction、learning}。

Stone 在开场明确表示，本讲要覆盖的是一个 broad view，而不是只看 2025 年流行的 ``agentic LLM''。因此课程的组织方式也很典型：先给概念和研究版图，再快速扫过实验线索，最后对两个代表性案例做 deeper dive（Slack 与 GT Sophy）。

\begin{importantbox}{本讲的核心价值}
这节课最关键的贡献不是新算法，而是把 Agent 研究中的三个长期主题重新拧到一起：
\begin{itemize}
    \item \textbf{Interaction}：单体与多体如何在动态环境中协作与对抗；
    \item \textbf{Learning}：从 RL 到 human feedback，如何在样本约束下学习可执行策略；
    \item \textbf{Embodiment}：策略是否真正能落在 physical world，而不是只在 text benchmark 上成立。
\end{itemize}
\end{importantbox}

\begin{knowledgebox}{为什么这讲对 LLM Agent 仍然重要}
尽管讲者大量举的是 robotics 与 autonomous driving 的例子，但这些问题与 LLM Agent 的工程现实高度同构：状态观测不完美、动作后果延迟显现、多人系统中的策略耦合、以及 ``高分但不可用'' 的评测失真。
\end{knowledgebox}

\subsection{本章小结}
这是一节 ``概念澄清 + 案例复盘'' 的课。它要求我们先把 Agent 的定义站稳，再去谈任何看似新的 Agentic 产品形态。

